\DocumentMetadata{
  pdfversion=2.0,
  pdfstandard=ua-2,
  lang=en-US,
  tagging=on,
  tagging-setup={math/setup=mathml-SE,tabsorder=structure}
}
\documentclass[11pt,letterpaper]{article}
\usepackage[margin=0.68in,includefoot]{geometry}
\usepackage{newcomputermodern}
\usepackage{amsmath,amscd,mathtools}
\usepackage{tikz,adjustbox}
\providecommand{\square}{\mdlgwhtsquare}
\usepackage[hidelinks]{hyperref}
\hypersetup{
  pdftitle={Analysis PhD Exam, August 2025},
  pdfauthor={University of Florida Department of Mathematics},
  pdfsubject={Graduate mathematics examination},
  pdfdisplaydoctitle=true,
  bookmarksopen=true
}
\setcounter{secnumdepth}{0}
\setlength{\parindent}{0pt}
\setlength{\parskip}{0.28em}
\setlength{\emergencystretch}{3em}
\setlength{\leftmargini}{2.15em}
\setlength{\leftmarginii}{2.65em}
\setlength{\leftmarginiii}{3em}
\pagestyle{empty}
\raggedbottom
\allowdisplaybreaks
\NewTaggingSocketPlug{block/list/label}{examproblem}{%
  \tagstructbegin{tag=Lbl}\tagstructend%
  \tagstructbegin{tag=\LBody}%
  \tagstructbegin{tag=\ProblemHeadingTag,title-o={\ProblemHeadingText}}%
  \tagmcbegin{tag=\ProblemHeadingTag,actualtext={\ProblemHeadingText}}%
  #2%
  \tagmcend%
  \tagstructend%
}
\newenvironment{examproblems}[2]{%
  \def\ProblemHeadingTag{#2}%
  \AssignTaggingSocketPlug{block/list/label}{examproblem}%
  \begin{enumerate}%
  \setcounter{enumi}{#1}%
  \renewcommand{\labelenumi}{\textbf{\theenumi.}}%
  \setlength{\topsep}{0.35em}%
  \setlength{\partopsep}{0pt}%
  \setlength{\itemsep}{0.48em}%
  \setlength{\parsep}{0.18em}%
}{\end{enumerate}}
\newenvironment{examsubparts}{%
  \AssignTaggingSocketPlug{block/list/label}{default}%
  \begin{enumerate}%
  \setlength{\topsep}{0.18em}%
  \setlength{\partopsep}{0pt}%
  \setlength{\itemsep}{0.16em}%
  \setlength{\parsep}{0.1em}%
}{\end{enumerate}}
\newenvironment{examinstructions}{%
  \par\begingroup\itshape
}{%
  \par\endgroup\vspace{0.18em}
}
\makeatletter
\renewcommand\subsection{\@startsection{subsection}{2}{0pt}%
  {-1.25ex plus -.25ex minus -.1ex}%
  {0.55ex plus .1ex}%
  {\normalfont\large\bfseries}}
\renewcommand{\@maketitle}{%
  \newpage\null\vskip -1em%
  {\footnotesize\bfseries
    \href{https://www.ufl.edu/}{University of Florida}, \href{https://math.ufl.edu/}{Department of Mathematics}\par}%
  \vskip -0.5em%
  \tagstructbegin{tag=H1,title={Analysis PhD Exam, August 2025}}%
  \tagpdfparaOff%
  \tagmcbegin{tag=H1}%
    {\LARGE\bfseries\href{https://gma.math.ufl.edu/exams/analysis-phd/}{Analysis PhD Exam}, August 2025\par}%
  \tagmcend%
  \tagpdfparaOn%
  \tagstructend%
  \vskip 0.25em%
  \tagmcbegin{artifact}%
    \hrule height 1pt%
  \tagmcend%
  \vskip 1em%
}
\makeatother
\title{Analysis PhD Exam, August 2025}
\author{}
\date{}
\begin{document}
\maketitle
\thispagestyle{empty}
\begin{examinstructions}
Write solutions in a neat and logical fashion, giving complete reasons for all steps. Do six of the eight problems.
\end{examinstructions}
\begin{examproblems}{0}{H2}
\def\ProblemHeadingText{Problem 1.}
\item
Do two. Short answer, you do not need to give a detailed proof, just a brief explanation.
\begin{examsubparts}
\item[\textnormal{(a)}]
If \(a_1,\ldots,a_n\) are non-negative numbers, must it be the case that
\[
\left(\sum_{j=1}^n a_j\right)^5 \leq n^4\sum_{j=1}^n a_j^5?
\]
\item[\textnormal{(b)}]
Define \(f\in L^1(\mathbb{R})\) by \(f(x)=e^{-x^4}\chi_{[1,\infty)}(x)\), where~\(\chi\) denotes indicator function. Is the Fourier transform of~\(f\) in \(L^1(\mathbb{R})\)?
\item[\textnormal{(c)}]
If \(g,h\in L^1(\mu)\), when are \(g\,d\mu\) and \(h\,d\mu\) mutually singular?
\end{examsubparts}
\def\ProblemHeadingText{Problem 2.}
\item
Suppose \((X,\mathcal{M})\) is a measurable space. Show, if \(E\in\mathcal{M}\otimes\mathcal{M}\), then
\[
\{x\in X:(x,x)\in E\}\in\mathcal{M}.
\]
\def\ProblemHeadingText{Problem 3.}
\item
This problem has two parts.
\begin{examsubparts}
\item[\textnormal{(a)}]
Let \(f_n(x)=x^{-n}\) for \(n\geq 2\). Does the numerical sequence
\[
c_n=\int_1^\infty e^{inx}f_n(x)\,dx
\]
 converge?
\item[\textnormal{(b)}]
Give an example, if possible, where strict inequality holds in Fatou’s Lemma.
\end{examsubparts}
\def\ProblemHeadingText{Problem 4.}
\item
Suppose~\(X\) is a Banach space, \(Y,Z\subseteq X\) are (closed) subspaces and \(Y\cap Z=\{0\}\). Prove or disprove: there is a \(C>0\) such that
\[
C(\lVert y\rVert+\lVert z\rVert)\leq\lVert y+z\rVert
\]
 for all \(y\in Y\) and \(z\in Z\).
\def\ProblemHeadingText{Problem 5.}
\item
Highlighting the key points, outline a proof of the following statement. If a sequence of functions converges in \(L^1\), then it has a subsequence that converges pointwise almost everywhere.
\def\ProblemHeadingText{Problem 6.}
\item
Suppose \((X,\mathcal{M},\mu)\) is a measure space, \(1\leq p,q\leq\infty\) and \(h:X\to\mathbb{C}\) is measurable. Show, if \(hf\in L^q(\mu)\) for each \(f\in L^p(\mu)\), then the map
\[
L^p(\mu)\ni f\mapsto hf\in L^q(\mu)
\]
 is continuous.
\def\ProblemHeadingText{Problem 7.}
\item
Define the Hardy–Littlewood maximal function of a Lebesgue measurable function \(f:\mathbb{R}\to\mathbb{R}\). State the Hardy–Littlewood maximal inequality. Show, if \(f\in L^1(\mathbb{R})\), then \(f^*\) (the Hardy–Littlewood maximal function of~\(f\)) is finite almost everywhere.
\def\ProblemHeadingText{Problem 8.}
\item
Do two. Short answer, you do not need to give a detailed proof, just a brief explanation.
\begin{examsubparts}
\item[\textnormal{(a)}]
Does there exist a bounded linear functional \(\lambda:L^\infty(\mathbb{R})\to\mathbb{C}\) for which there does not exist a \(g\in L^1(\mathbb{R})\) such that \(\lambda(f)=\int fg\,dx\)?
\item[\textnormal{(b)}]
Does the norm in \(\ell^1(\mathbb{N})\) arise from an inner product?
\item[\textnormal{(c)}]
Can a sigma algebra be bijective to~\(\mathbb{N}\)?
\end{examsubparts}
\end{examproblems}
\end{document}
